\begin{figure}[htbp]
	\centering		
	\begin{tikzpicture}[>=Stealth, line width=0.8pt]
		
		% ═══════════════════════════════════════════════════════════
		% PANNELLO (a) — Simbolo del carrello
		% ═══════════════════════════════════════════════════════════
		\begin{scope}[xshift=0cm]
			
			% ── Parametri del simbolo (convenzioni del libro) ──
			\def\xA{0.8}     % ascissa di A
			\def\yA{0.0}     % ordinata di A
			\def\tw{0.30}    % semi-larghezza triangolo
			\def\thR{0.40}   % altezza triangolo carrello
			\def\rn{0.10}    % raggio semicerchio (anellino a metà)
			\def\bw{0.18}    % semi-distanza ruote
			\def\gw{0.45}    % semi-larghezza barre di terra
			
			% ── Asse del carrello: retta verticale grigia continua ──
			\draw[gray!60, line width=0.6pt] 
			(\xA, \yA-\thR-0.5) -- (\xA, 3.85);
			
			% ── Etichetta del luogo del CR ──
			\node[gray, anchor=west, font=\footnotesize] 
			at (\xA+0.08, 2.50) {$C \in$ asse del carrello};
			
			
			% ── Retta-luogo del CR: due tratti, sopra il corpo e sotto il suolo ──
			\draw[dashed, gray!60, line width=0.6pt] 
			(\xA, 2.05) -- (\xA, 2.85);
			\draw[ gray!60, line width=0.6pt] 
			(\xA, \yA-\thR) -- (\xA, \yA-\thR-0.15);
			
			% ── Corpo di fantasia (linea, senza fill) ──
			\draw[line width=1pt] 
			plot[smooth cycle, tension=0.7] coordinates 
			{(0.8, 0.0) (1.6, 0.15) (2.2, 0.7) (2.0, 1.7) 
				(1.0, 2.0) (-0.1, 1.7) (-0.6, 1.0) (-0.3, 0.3)};
			\node at (1.5, 1.1) {$\mathcal{C}$};
			
			% ── Spostamento ammesso: freccia grigia perpendicolare ad a ──
			%\draw[->, gray, line width=0.9pt] 	(\xA+0.20, \yA) -- (\xA+1.60, \yA);
			\node[above=5pt] at (\xA+2.50, \yA+0.5) 	{\footnotesize $\ub(A) \cdot \ab = 0$};
			
			% ── Simbolo del carrello: apice del triangolo IN A ──
			
			% Triangolo (apice in A, base in basso)
			\draw (\xA, \yA) -- ++(-\tw, -\thR) -- ++(2*\tw, 0) -- cycle;
			
			% Prima barra (alla base del triangolo)
			\draw (\xA-\gw, \yA-\thR) -- (\xA+\gw, \yA-\thR);
			
			% Due ruote (cerchietti vuoti)
			\draw (\xA-\bw, \yA-\thR-0.13) circle (3pt);
			\draw (\xA+\bw, \yA-\thR-0.13) circle (3pt);
			
			% Seconda barra (suolo)
			\draw (\xA-\gw, \yA-\thR-0.26) -- (\xA+\gw, \yA-\thR-0.26);
			
			% Tratteggio del suolo
			\foreach \x in {-0.36,-0.18,0,0.18,0.36}
			\draw[thin] (\xA+\x, \yA-\thR-0.26) -- ++(-0.10, -0.13);
			
			% Semicerchio (anellino a metà) in A: rotazione consentita
			\draw[fill=white] 
			(\xA-\rn, \yA) arc(180:360:\rn) -- cycle;
			
			% Punto A (ridisegnato dopo il semicerchio così resta visibile)
			%\fill (\xA, \yA) circle (1.5pt);
			\node[above left=-1pt] at (\xA-0.02, \yA+0.05) {$A$};
			
			% ── Asse a sopra il corpo ──
			\draw[->, thick] (\xA, 2.10) -- (\xA, 2.85);
			\node[right=2pt] at (\xA, 2.85) {$\ab$};
			
			% ── Etichetta verticale "asse del carrello" ──
			\node[gray, rotate=90, anchor=south, font=\footnotesize] 
			at (\xA-0.12, 2.45) {asse del carrello};
			
			% ── Etichetta pannello ──
			\node[anchor=north] at (\xA+2, \yA-\thR) {$(a)$};
			
		\end{scope}
		
		
		% ═══════════════════════════════════════════════════════════
		% PANNELLO (b) — biella (realizzazione meccanica)
		% ═══════════════════════════════════════════════════════════
		\begin{scope}[xshift=5.5cm]
			
			% ── Parametri ──
			\def\xA{0.8}     % ascissa di A (e del perno al suolo)
			\def\yA{1.5}     % ordinata di A
			\def\yG{0.0}     % ordinata del perno al suolo
			\def\tw{0.30}    % semi-larghezza triangolo
			\def\thA{0.40}   % altezza triangolo cerniera al suolo
			\def\rn{0.10}    % raggio semicerchio in A
			\def\gw{0.50}    % semi-larghezza barra di terra
			
			% ── Asse della biella: retta verticale dashed continua ──
			%    (verrà parzialmente coperta da asta e triangolo)
			\draw[gray!60, line width=0.6pt] 
			(\xA, \yG-\thA-0.5) -- (\xA, 4.50);
			
			% ── Corpo di fantasia (alzato, senza fill) ──
			\draw[line width=1pt] 
			plot[smooth cycle, tension=0.7] coordinates 
			{(0.8, 1.5) (1.6, 1.65) (2.2, 2.2) (2.0, 3.2) 
				(1.0, 3.5) (-0.1, 3.2) (-0.6, 2.5) (-0.3, 1.8)};
			\node at (0.95, 2.6) {$\mathcal{C}$};
			
			% ── Cinematica finita: arco ──
			\draw[gray, line width=0.9pt] 
			(\xA, \yG) ++(60:1.5) 
			arc[start angle=60, end angle=120, radius=1.5];
			\node[gray, anchor=west, font=\footnotesize] 
			at (1.6, 0.85) {cinematica finita};
			
			% ── Cinematica infinitesima: tangente ──
			\draw[gray, dashed, line width=0.9pt] 
			(-0.7, \yA) -- (2.4, \yA);
			\node[gray, anchor=west, font=\footnotesize] 
			at (2.2, \yA+0.35) {cinematica infinitesima};
			
			% ── Asta della biella (copre il tratto centrale dell'asse) ──
			\draw[line width=0.8pt] (\xA, \yA) -- (\xA, \yG);
			
			% ── Cerniera al suolo: triangolo + barra + tratteggio ──
			\draw (\xA, \yG) -- ++(-\tw, -\thA) -- ++(2*\tw, 0) -- cycle;
			\draw (\xA-\gw, \yG-\thA) -- (\xA+\gw, \yG-\thA);
			\foreach \x in {-0.36,-0.18,0,0.18,0.36}
			\draw[thin] (\xA+\x, \yG-\thA) -- ++(-0.10, -0.13);
			
			% ── Anellini (in quest'ordine, dopo asta e triangolo) ──
			% Sommità della biella (in A): semicerchio bianco, concavità in basso
			\draw[fill=white] 
			(\xA-\rn, \yA) arc(180:360:\rn) -- cycle;
			% Cerniera al suolo: perno pieno
			\draw[fill=white] (\xA, \yG) circle (3.5pt);
			
			% ── Punto A (ridisegnato dopo il semicerchio) ──
			%\fill (\xA, \yA) circle (1.5pt);
			\node[above left=-1pt] at (\xA-0.02, \yA+0.05) {$A$};
			
			% ── Asse a sopra il corpo ──
			\draw[->, thick] (\xA, 3.60) -- (\xA, 4.30);
			\node[right=2pt] at (\xA, 4.30) {$\ab$};
			
			% ── Etichetta del luogo del CR ──
			\node[gray, anchor=west, font=\footnotesize] 
			at (\xA+0.08, 4.00) {$C \in$ asse della biella};
			
			% ── Etichetta pannello ──
			\node[anchor=north] at (\xA+2, \yG-\thA) {$(b)$};
			
		\end{scope}
		
	\end{tikzpicture}
	\caption{Vincolo semplice di carrello: simbolo convenzionale (a)
		e realizzazione meccanica tramite biella (b). L'equazione
		di vincolo $\ub(A)\cdot\ab = 0$ annulla la componente di
		spostamento di $A$ lungo l'asse $\ab$; le componenti libere
		--- traslazione lungo $\ab^\perp$ e rotazione --- non sono
		rappresentate. Per la biella, la cinematica finita è un
		arco di circonferenza centrato nel perno al suolo;
		nell'ipotesi di spostamenti infinitesimi essa coincide con
		la tangente in $A$, e le prestazioni cinematiche della
		biella si riducono a quelle del carrello. In entrambi i
		casi il \CRa\ giace sulla retta per $A$ in direzione $\ab$.}
	\label{fig:carrello}
\end{figure}
